\section{Introduction}
\label{sec:intro}

Clinical records are encrypted at rest because regulation requires it. The HIPAA Security Rule lists encryption of electronic protected health information among its technical safeguards~\cite{hipaa-security-rule}; the General Data Protection Regulation names health data a special category and requires security appropriate to the risk~\cite{gdpr}. The specification that results from these obligations, in procurement documents and in system designs, is almost always a cipher name and a key length. ``AES-256 at rest'' is a sentence that satisfies an auditor.

It is also a sentence about bytes, not about records. A cipher with a random nonce guarantees that ciphertext reveals nothing about plaintext, and an authentication tag guarantees that the bytes were produced under the key. Neither says anything about \emph{whose} record a ciphertext belongs to, in \emph{which} tenant, in \emph{which} field. That binding lives entirely in the surrounding database: in the row key, the table, the column name. An adversary who can write to the database can rewrite all three. An insider with write access, a database administrator, a compromised service account, or a migration script with a bug, can copy the encrypted note from one patient's row into another patient's row. The application later reads the second row, decrypts under the key it holds, and the tag verifies, because the bytes are genuine. A clinician sees one patient's data under another patient's name, and nothing in the system has failed. This is not a weakness in AES. It is a property that the deployment never asked for.

The mechanism to ask for it has existed for two decades. Authenticated encryption with associated data~\cite{rogaway2002aead} lets a caller bind a ciphertext to arbitrary public data that must be presented again at decryption; cloud key services expose the same idea as an encryption context~\cite{aws-kms-context}. It is rarely used in storage encryption for clinical data, and where it is used, the choice of what to bind, and how to encode it so that two contexts can never collide, is left to each application. Section~\ref{sec:evaluation} confirms that associated data alone stops every substitution we try. What it does not do is limit what a single leaked key exposes, or make a retired key stop working; those are the gaps between an associated-data deployment and the construction of this paper, and they are measured rather than asserted.

A separate literature has taken a different route to ``more than AES''. A body of work on hybrid encryption for health records places a keyed byte-level transform, an affine map, a substitution, a chaotic sequence, a polynomial, in front of a standard cipher and reports added security. Section~\ref{sec:pretransform} examines what such a layer contributes: nothing if its key travels with the ciphertext, as it must when the receiver does not already share it; no diffusion in any case; and, for the most common instance, a keyspace of $2^{\KeyspaceLog}$, shown by exhaustive check. The engineering effort those layers absorb is better spent on binding, and this paper spends it there.

The question a reader should put to this paper is whether a specification made of known parts is a contribution. Section~\ref{sec:design} says plainly that nothing is new at the primitive level, and Section~\ref{sec:discussion} shows how much of the same discipline the AWS Database Encryption SDK already applies to DynamoDB. What is claimed is narrower and, we think, more useful than a new primitive: a statement of which facts a clinical ciphertext must be bound to and an encoding under which that binding cannot be got wrong, a key hierarchy whose consequences for a leaked key and for its containment are measured rather than asserted, and an evaluation that compares the result against the strongest thing the AEAD interface offers without it, not against the weakest. The pre-transform analysis of Section~\ref{sec:pretransform} is the negative counterpart: a measurement of what a widely proposed alternative does not buy.

\paragraph{Contributions}
\begin{enumerate}
  \item A precise account of what a byte-level pre-transform layer in front of a standard cipher can and cannot add, with an exhaustive check of the most common instance and a fair measurement of its cost (Section~\ref{sec:pretransform}).
  \item CPDPM, a context-bound authenticated-encryption construction for clinical records: an injective encoding of tenant, patient, resource type and field path; per-context keys derived from a versioned master key; the same context bound as associated data; a key commitment; a single failure mode; an encrypted manifest binding the set of encrypted fields in a FHIR resource; and a streaming mode. Security reduces to the standard AEAD notions (Section~\ref{sec:design}).
  \item An empirical comparison along a ladder of four schemes, shared-key AES-CBC+HMAC, per-tenant AES-GCM, per-tenant AES-GCM with the context as associated data, and CPDPM, showing what each step buys: associated data alone stops substitution across patients, tenants and fields; the per-context key reduces the exposure of one leaked key from every record of a tenant to one field of one patient, and containing that leak to re-encrypting one record with the fewest secrets at rest of any scheme; versioning rejects stale ciphertext after key retirement. And a measurement of cost against raw AES-GCM, Fernet and Tink on two architectures (Section~\ref{sec:evaluation}).
  \item An open artifact: the implementation, a test suite with property-based fuzzing of every parser, known-answer vectors, and every experiment and figure, reproducible with one command.
\end{enumerate}

Section~\ref{sec:background} places the work against authenticated encryption, key commitment, encryption context in practice, and the hybrid-encryption literature. Section~\ref{sec:discussion} states what the construction does not do.
